\sscp{The form of the reconstruction as **mantəR ‘cuscus’}{14-04-03}
We accept that local borrowing might account for some of the
Sulawesi data, but the consistency of form
and meaning in two separate branches of the \tsc{\ili{Muna}-Buton}
subgroup \citep{vandenBerg2003} mean that borrowing is unlikely
to explain all of them.\footnote{
		A reviewer points out that the
		fronting affect of *R is only regular in Munan languages (\ili{Kaimbulawa},
		\ili{Pancana}, \ili{Kioko}, \ili{Kambowa}, and \ili{Busoa} in \trf{14-07a}.
		See \trf{14-09} for more data on fronting).
		Given this, \ili{Cia-Cia} and \ili{Kaisabu} \it{mante} ‘bear cuscus’ are
		likely borrowings from a Munan language.}
The \it{mante} forms
in two separate branches of the \tsc{\ili{Muna}-Buton} subgroup in
Sulawesi raise several natural questions:

\begin{itemize}
	\item	Could the Sulawesi forms be shown to be related to putative
				reflexes of **mans(aə)r reflexes,
				rather than just similar?
	\item	Given the form \it{mante} in several languages,
				and the paucity of forms with medial \it{s} throughout
				Wallacea, is **mans(aə)r actually the best reconstruction to account for
				the entirety of the data relating to ‘cuscus’ in the target region?\footnote{
						While we are not asserting a common subgroup, or parent language,
						nevertheless it is worth exploring to see if the diversity of
						the surface forms can be accounted for a little better.}
				Specifically:
					\begin{itemize}
						\item	Is medial *ns the best way to account for the medial consonants
									in the greatest number
									and distribution of languages,
									or might *nt or *nd provide an equally good or,
									better account of the wider data?
						\item	Given reconstructions of both **mansar
									and **mansər, can we gain better insights into the vowel of the
									ultimate syllable? \item	Since reflexes of PMP *r are so rare (found in only a few
									etyma for which cognates are often hard to find),\footnote{
											\citet{Wolff2010}	is one who does not accept \tsc{PAn}/PMP *r
											as a proto-phoneme. See \srf{01-07}.}
									could the final consonant be better accounted
									for by *R?
					\end{itemize}
\end{itemize}